\section{Critical-roof and fold audit: normalization and ordered closure}
\label{sec:cr-audit}
The new identities have different kinds of uniformity. The
clearance-to-endpoint determinant is an identity for each finite
physical itinerary and every marked flight. The fold and finite-type
bounds are uniform over Borel selectors on a specified chart, with
specified derivative constants. The critical coarea error is uniform
over selectors but depends on the chart's third-derivative and
weight-modulus constants. The critical coefficient sum is over all
regular physical words at fixed collision count, relative to the
retained uniform pressure input; it is unscaled. None of these
qualifiers can be replaced by another.

\subsection{Two distinct uses of reversibility}
At a normal endpoint, the actual reflected ray reverses its incoming
velocity. This proves the reversing-periodic-orbit assertion without
adding a reflection at an unvisited mark. In collision coordinates,
time reversal also gives the derivative identity
$\mathcal M=JB^{-1}JB$. The first assertion concerns actual points
and their section tests; the second concerns the tangent map on the
regular domain. At an unvisited tangency, only the one-sided reduced
action coefficient is used. A smooth extension of that action does
not define the physical map through the forbidden side.

Even at a regular reversing orbit the original section need not be
reversal invariant. The occupation identity is
$A_{2m}=A_m+A_m^{\mathcal I}$ when both original endpoints are in
the section, not automatically $2n$. The coefficient
$\mathfrak J$ describes nearby collision area pushed to the roof.
It is not a conditional law on a measure-zero periodic orbit.

\subsection{The exact small-scale error and its range}
The radius $r_0$ in Theorem~\ref{thm:cr-critical-density} is a local
endpoint-chart radius. The height increment is $s=t-t_x$, and its
allowed range is $0<s<r_0^2/3$. Neither $r_0$ nor $s$ replaces the
original microscopic observation window. In particular the explicit
error constant
\[
 2\pi\sqrt3\{2L_w+Kw(0)\}
\]
has not been shown summable after multiplying the full physical
source by $m^2$. Decreasing $r_0$ separately for each orbit does not
establish that summability. A general Borel selector retains the
exact fraction $a_E(s)$; only a verified finite tangent-sector
structure justifies a limiting angular fraction.

The local fold exponent $1/2$ is produced by integrating a square-root
conditional singularity across a transverse clearance direction.
A higher finite type has its stated exponent $1/\ell$ and derivative
budget. Full critical roofs have a different two-dimensional local
coarea law; the fold theorem is not applied there when its clearance
transversality fails. A countable disjoint source-chart covering
preserves positivity and exact multiplicity, but does not itself
bound the sum of the chart constants.

\subsection{The unchanged closure statement}
All original positive coefficients and the complete nongrazing pool
are still those of Section~\ref{sec:v112-source}. The remaining
statement is \eqref{eq:v112-remaining}, with the full physical sum
first, then the exact-label and radius suprema and the original roof
essential supremum, then the collision limsup, and the outer limit
last. A successful continuation must control clearance and critical
charts, positive one-sided pairing traces, the low-incidence source,
and the actual available comparison capacity in that topology.
The new calculations neither subtract these unpaid classes nor
introduce a new name for the shortage.

The accompanying programs check the signed generating derivative,
stability determinant, normalization, fold constants, selector error
and actual horizontal period-two physical rays. The latter are not
asserted to lie in the large-collision central roof window. Exact
rational checks of their optical Schur complements and finite
ray-derivative comparisons test independent implementations; they are
not substitutes for the written continuous proofs or the unresolved
uniform dynamical estimates.
